\documentclass[a4paper,12pt]{ujarticle}
\usepackage[top=25mm,bottom=25mm,left=25mm,right=25mm]{geometry}
\title{週間報告}
\author{金武俊佑}
\date{2025年3月16日}

\usepackage{otf}
\usepackage{url}
\usepackage{tabularx}
\usepackage{comment}
\usepackage[dvipdfmx]{graphicx}
\usepackage[dvipdfmx]{color}
\usepackage{listings}
\usepackage{pdfpages}
\usepackage{multirow}
\usepackage{float}

\newcommand{\thc}[1]{\multicolumn{1}{c|}{#1}}

\begin{document}
\maketitle

\section{今週の報告}\label{sec:thisweek}
今週の報告を以下に記す。
\subsection{就活}\label{sec:job}
トランスコスモスのESを書いた。時間があるときにキャリアセンターでESの確認と面接の対策を行いたい。
\subsection{その他}\label{sec:other}
ペルソナ5Rを買った。3月になって自由な時間が無くなってきているので、時間があるときに少しずつ進めていきたい。
\section{来週の予定}\label{sec:nextweek}
来週の活動予定を以下に記す。
\begin{itemize}
  \item 03/21(土) 10:30 miniPBLに参加する。
\end{itemize}
\section{ 卒論の候補}\label{sec:soturon}
卒論の候補を以下に記す。

\subsection{AI会話練習・面接練習支援システムの開発}

\subsubsection{内容}
本テーマでは、孤独感やコミュニケーション不安のある学生を対象に、AIと雑談や面接練習ができる対話型システムの開発を目指す。利用者は雑談モードや面接練習モードを選択し、AIと会話する。対話後には、話の長さ、結論から話せているか、具体性、質問への答え方などについてフィードバックを返す機能を持たせることを考えている。
このシステムにより、会話や面接に対する心理的ハードルを下げ、不安軽減につながることを期待している。実装面では、LLM APIを利用することで対話機能は比較的実現しやすく、Webアプリケーションとして構築可能だと考えている。評価については、使用前後のアンケートにより、会話への不安や面接への自信の変化を調べることを想定している。


\subsubsection{実現可能かどうか}
LLM APIを利用することで対話機能の実装は可能であり、Webアプリケーションとして構築しやすいと考える。対話履歴の保存や簡易的なフィードバック機能も実装可能である。評価については、使用前後のアンケートによって会話不安や面接への自信の変化を調べる方法が考えられる。

\subsection{ネット記事の共同閲覧体験を模倣するAI対話支援システムの開発}

\subsubsection{内容}

本テーマでは、利用者がネット記事を読む際に、AIが一緒に記事を読んでいる相手のように感想や質問、別の視点を返すシステムの開発を目指す。単なる要約ツールではなく、「友達と一緒に記事を読んで話す」ような体験をAIによって再現することを考えている。
このシステムにより、一人で記事を読むときでも内容理解を深めたり、感想を言語化しやすくしたりできると考えられる。また、記事閲覧時の孤独感の軽減にもつながる可能性がある。実装面では、記事本文を取得または入力し、その内容をLLM APIに渡して要約、感想、質問生成などを行うWebアプリとして構築できると考えている。評価については、AIなしで記事を読んだ場合とAIありで読んだ場合を比較し、理解のしやすさや感想の書きやすさ、閲覧体験の楽しさなどをアンケートで調べることを想定している。


\subsubsection{実現可能かどうか}
記事の本文を取得または入力し、その内容をLLM APIに渡すことで、要約、感想、質問生成などの対話機能は実装可能である。まずはURLまたは本文入力型のWebアプリとして構築すれば、技術的にも現実的である。評価については、AIなしで読んだ場合とAIありで読んだ場合を比較し、理解しやすさや感想の書きやすさ、閲覧体験の楽しさをアンケートで測定することが考えられる。
\subsection{AIを用いた気分整理・思考整理支援システムの開発}

\subsubsection{内容}
本テーマでは、利用者が不安や悩み、考えごとを文章で入力すると、AIが内容を整理し、気持ちや課題を分かりやすく言語化する支援システムの開発を目指す。入力内容を「不安」「迷い」「やるべきこと」などに分けて整理し、自分の状態を振り返りやすくすることを考えている。
このシステムはだ、単なる相談システムではなく、自分の気持ちや考えを整理するための補助を行うことを目的とする。実装面では、文章入力、履歴保存、AIによる整理結果の表示といった構成で、Webシステムとして実現可能と考えている。評価については、使用前後で「頭の中を整理しやすくなったか」「自分の考えを言葉にしやすかったか」などをアンケートで調べる方法を想定している。
\subsubsection{実現可能かどうか}
文章入力、履歴保存、AIによる整理結果の表示といった構成であれば、Webシステムとして実装可能である。LLM APIを利用すれば、入力内容の分類や要約、整理支援の応答も比較的実現しやすい。評価については、使用前後で「頭の中を整理しやすくなったか」「自分の考えを言葉にしやすかったか」といった観点からアンケートを行うことが考えられる。

\subsection*{まとめ}
以上の3つの候補の中から、社会的意義、実装可能性、自分の興味関心を踏まえて、最終的な卒業研究テーマを決定したいと考えている。



\section{スケジュール}\label{sec:schedule}
後期のスケジュールを表\ref{tab:studyschedule}に示す。

\begin{table}[H]
  \caption{後期のスケジュール}\label{tab:studyschedule}
  \centering
  \begin{tabular}{|p{4mm}||p{30mm}|p{30mm}|p{30mm}|p{30mm}|p{30mm}|}\hline
        & \thc{月} & \thc{火} & \thc{水} & \thc{木} & \thc{金} \\ \hline\hline
     1  &     & AI導入    &情報セキュリティ     & キャリア形成 & テクノアート  \\ \hline
     2  &     & 知的財産権 &  & 現代の政治 & テクノアート \\ \hline
     3  &     & 経営情報学 & 情報科学演習Ⅱ &  &  \\ \hline
     4  &  &  & &  & \\ \hline
  \end{tabular}
\end{table}

\end{document}
